\section{Step 20 Compare a round boundary and the complete block result}\label{sha-s20}

First pause after the working-vector replacement and compare the eight words in the worked example. Then let the completed function return and compare the feed-forward state and preserved registers. The two observations answer different questions: how a round changes working storage and what the public function returns to its caller.

\subsection{Observe and compare}
Use recipes \texttt{s20-1}, \texttt{s20-2}. Prepare each in a separate source image from the current cumulative program. An internal checkpoint is inspected before its temporary or normal return; otherwise pause at \texttt{main\_done}.

Use the completed entries available at this point; keep their inputs and observation boundary explicit. The surrounding explanation identifies the state being compared.

\subsection{Expected observation}
The first round matches the table; the returned state matches compression and preserved registers have their original full values.

\Needspace{9\baselineskip}
\subsection{Check your understanding}
Should \texttt{s0} still hold working a after the function returns?

\begin{itemize}
\item \textbf{A} Yes, saved registers are the digest output.
\item \textbf{B} No; the public function restores the caller's \texttt{s0}.
\item \textbf{C} Only when a is negative.
\end{itemize}
